\documentclass[10pt,a4paper]{article}
\usepackage{amsmath,amssymb,amsthm,graphicx}
% Card-local minimal preamble (journal unit sigma-2025-020-non-integrability-of-the-sasano-system-of-ty, P2-ahl16 precedent).
%% Verbatim from the pinned source: the author theorem environments and macro block,
%% so every macro and environment of the counted spans is available.
% Omitted and why: the documentclass line (the journal class sigma is not on the host;
%% the assembler supplies the class) and page-geometry/font blocks (they fight the
%% forced xelatex engine).
\usepackage{amsmath,amssymb,amsthm,mathtools}

\newcommand{\RS}[1]{\textcolor{red}{#1}}
\numberwithin{equation}{section}
\newtheorem{Theorem}{Theorem}[section]
\newtheorem*{Theorem*}{Theorem}
\newtheorem{Corollary}[Theorem]{Corollary}
\newtheorem{Lemma}[Theorem]{Lemma}
\newtheorem{Proposition}[Theorem]{Proposition}
 { \theoremstyle{definition}
\newtheorem{Definition}[Theorem]{Definition}
\newtheorem{Note}[Theorem]{Note}
\newtheorem{Example}[Theorem]{Example}
\newtheorem{Remark}[Theorem]{Remark} }
\usepackage{cleveref}
\crefname{originaltheorem}{original theorem}{original theorems}
\Crefname{originaltheorem}{Original theorem}{Original theorems}
\makeatletter
\newlabel{Mi}{{2.11}{source}{}{sourcepointer-1}{}}
\newlabel{Mi@cref}{{[originaltheorem][0][]2.11}{[1][0][]source}{}{}{}}
\makeatother
\crefname{originaltheorem}{original theorem}{original theorems}
\Crefname{originaltheorem}{Original theorem}{Original theorems}
\makeatletter
\newlabel{reg}{{2.4}{source}{}{sourcepointer-2}{}}
\newlabel{reg@cref}{{[originaltheorem][0][]2.4}{[1][0][]source}{}{}{}}
\makeatother
\crefname{originalequation}{original equation}{original equations}
\Crefname{originalequation}{Original equation}{Original equations}
\makeatletter
\newlabel{fve}{{3.2}{source}{}{sourcepointer-3}{}}
\newlabel{fve@cref}{{[originalequation][0][]3.2}{[1][0][]source}{}{}{}}
\makeatother
\title{The connected component of the differential Galois group of the first normal variational equation is Abelian}
\author{Stoyanova, Tsvetana}
\date{doi:10.3842/SIGMA.2025.020; CC BY 4.0}
\begin{document}
\maketitle
\setcounter{section}{3}
\setcounter{Theorem}{0}
\noindent\textit{Editorial note.} All delivered mathematical text is the author's own, copied from the pinned journal source file sources/sigma-2025-020-non-integrability-of-the-sasano-system-of-ty/pinned-original-source.tex (Non-Integrability of the Sasano System of Type D5(1) and Stokes Phenomena); no statement or proof step was rewritten, reordered or removed.\par\medskip
\noindent\textit{Editorial note.} The counted proof is the author's complete argument: All chain roots are named external theorems of the same article (Mitschi (2.11), Schlesinger (2.4), Ramis (2.10) as the implicit frame for the no-Stokes step, which is an author observation). The source label of this theorem is commented out in the pinned source; the census label ve1 is retained as identifier. The source typo e\textasciicircum{}\{2 pi i alpha\} matrix entries at lines 507-511 render as written (one zero entry dropped per row, verbatim carry). External navigation stubs: Mi (2.11), reg (2.4), fve (3.2).\par\medskip
\noindent\textit{Editorial note.} The label ve1 names the counted statement, and the proof environment immediately following it in the source is delivered verbatim as its proof.\par\medskip
\noindent\textit{Editorial note.} The counted statement is carried by the author's own Theorem environment in the article source.\par\medskip
\noindent\textit{Editorial note.} Journal-unit assembly: card-local preamble (author theorem environments and macros verbatim), article class, no bibliography in the counted span (the span carries no \textbackslash{}cite commands).\par\medskip
\noindent\textit{Editorial scope.} the counted statement and its own proof environment only\par\medskip
\noindent\textit{Allowed standard prerequisites (editorial).}\begin{itemize}
\item Mitschi's theorem (Theorem Mi, printed (2.11)): the differential Galois group is topologically generated by the local differential Galois groups at the singular points.
\item Schlesinger's theorem (Theorem reg, printed (2.4)): at a regular singular point the monodromy group is Zariski dense in the local differential Galois group.
\item The NVE\_1 system (fve) (printed (3.2)): the explicit four-dimensional first-order system along the particular solution.
\item Standard regular-singularity facts: a fundamental solution has the form P(t) t\textasciicircum{}A with P holomorphic and A constant, and the monodromy is e\textasciicircum{}\{2 pi i A\}; the formal monodromy around t = infinity is the inverse of the monodromy around t = 0.
\item Ramis's theorem (Theorem Ra, printed (2.10)) as the frame for the no-Stokes case: when no non-trivial Stokes phenomena are observed, the local Galois group is generated topologically by the formal monodromy and the exponential torus.
\end{itemize}\medskip
\ifdefined\hypertarget\hypertarget{sourcepointer-1}{}\fi
\noindent\textit{Source reference.} Original theorem 2.11: sources/sigma-2025-020-non-integrability-of-the-sasano-system-of-ty/pinned-original-source.tex, source lines 406--406. Navigation pointer invoked at line(s) 485: the counted text invokes that named theorem (Mitschi) of the same article as an input, which is not retained in this card, so the reference is declared as a pointer carrying the number the article itself prints for it (the section and the place of the statement inside that section, derived by the census v3.5 rule for the shared section-scoped theorem counter); the cited theorem is an allowed input and none of it is reproduced inside the span.\par\medskip
\ifdefined\hypertarget\hypertarget{sourcepointer-2}{}\fi
\noindent\textit{Source reference.} Original theorem 2.4: sources/sigma-2025-020-non-integrability-of-the-sasano-system-of-ty/pinned-original-source.tex, source lines 283--283. Navigation pointer invoked at line(s) 503: the counted text invokes that named theorem (Schlesinger) of the same article as an input, which is not retained in this card, so the reference is declared as a pointer carrying the number the article itself prints for it (the section and the place of the statement inside that section, derived by the census v3.5 rule for the shared section-scoped theorem counter); the cited theorem is an allowed input and none of it is reproduced inside the span.\par\medskip
\ifdefined\hypertarget\hypertarget{sourcepointer-3}{}\fi
\noindent\textit{Source reference.} Original equation 3.2: sources/sigma-2025-020-non-integrability-of-the-sasano-system-of-ty/pinned-original-source.tex, source lines 459--459. Navigation pointer invoked at line(s) 475: the counted text cites that numbered equation of the same article, which is not retained in this card, so the reference is declared as a pointer carrying the number the article itself prints for it (the section and the place of the equation inside that section, derived by the census v3.5 rule for the section-scoped equation counter); the cited equation is an allowed input and none of it is reproduced inside the span.\par\medskip
 \begin{Theorem}%\label{ve1}
 The connected component $(G_1)^0$ of the unit element of the differential Galois group $G_1$ of the $({\rm NVE})_1$ is Abelian.
 \end{Theorem}

 \begin{proof}

 The $({\rm NVE})_1$ have two singular points over $\mathbb{C}\mathbb{P}^1$: the points $t=0$ and $t=\infty$. The~origin is a regular singularity while $t=\infty$ is an irregular singularity.
 The system~\eqref{fve} admits a~fundamental matrix solution $\Phi(t)$ in the form
 \[
 \Phi(t)=\begin{pmatrix}
 {\rm e}^{-t} t^{\alpha-1} & 0 & -\alpha {\rm e}^t t^{-\alpha-1} & -\alpha {\rm e}^t t^{-\alpha-1}\\[0.1ex]
 0 & {\rm e}^{-t} t^{\alpha-1} & -\alpha {\rm e}^t t^{-\alpha-1} & -\alpha {\rm e}^t t^{-\alpha-1}\\[0.1ex]
 0 & 0 & {\rm e}^t t^{-\alpha+1} & 0\\[0.1ex]
 0 & 0 & 0 & {\rm e}^t t^{-\alpha+1}
 \end{pmatrix}.
 \]
 We will compute the Galois group $G_1$ of the $({\rm NVE})_1$ with respect to this fundamental matrix solution.
 In this case, from Theorem \ref{Mi}, it follows that
 the differential Galois group $G_1$ of the system~\eqref{fve} is generated topologically by the local Galois groups $G_0$ and $G_{\infty}$, corresponding to the singularities $t=0$ and $t=\infty$, respectively.

 Let us first determine the group $G_0$.
 In a neighborhood of the origin, the above solution $\Phi(t)$ is written as
 \[
 \Phi(t)=P(t) t^A ,
 \]
 where $P(t)$ is the holomorphic matrix
 \[
 P(t)=\begin{pmatrix}
 {\rm e}^{-t} & 0 & -\alpha {\rm e}^t & -\alpha {\rm e}^t\\
 0 & {\rm e}^{-t} & - \alpha {\rm e}^t & -\alpha {\rm e}^t \\
 0 & 0 & t^2 {\rm e}^t & 0\\
 0 & 0 & 0 & t^2 {\rm e}^t
 \end{pmatrix} .
 \]
 For the constant matrix $A$ we have that $A={\rm diag} (\alpha-1, \alpha-1, -\alpha-1, -\alpha-1)$.
 From Theorem~\ref{reg}, it follows that the Galois group $G_0$ over $\mathbb{C}(t)$ is generated topologically by the monodromy matrix~$M_0$ around the origin.
 With respect to the fundamental matrix solution $\Phi(t)$, we obtain
 \[
 M_0={\rm e}^{2 \pi {\rm i} A}=\begin{pmatrix}
 {\rm e}^{2 \pi {\rm i} \alpha} & 0 & 0 & 0\\
 0 & {\rm e}^{2 \pi {\rm i} \alpha} & 0\\
 0 & 0 & {\rm e}^{-2 \pi {\rm i} \alpha} & 0\\
 0 & 0 & 0 & {\rm e}^{-2 \pi {\rm i} \alpha}
 \end{pmatrix} .
 \]
 When $\alpha\in\mathbb{Q}$ but $\alpha\notin\mathbb{Z}$, the group generated by $M_0$ is not connected but it is a finite and cyclic group. In this case,
 \[
 G_0=\left\{\begin{pmatrix}
 \nu^{-1} & 0 & 0 & 0\\
 0 & \nu^{-1} & 0 & 0\\
 0 & 0 & \nu & 0\\
 0 & 0 & 0 & \nu
 \end{pmatrix}, \ \nu \ \text{is a root of unity} \right\}, \qquad (G_0)^0=\{I_4\} ,
 \]
 where $I_4$ is the identity matrix. When $\alpha\in\mathbb{Z}$, we have that $G_0=(G_0)^0=\{I_4\}$. When $\alpha\notin\mathbb{Q}$, the group generated by $M_0$ is a connected group
 and
 \[
 G_0=(G_0)^0=\left\{ \begin{pmatrix}
 \nu^{-1} & 0 & 0 & 0\\
 0 & \nu^{-1} & 0 & 0\\
 0 & 0 & \nu & 0\\
 0 & 0 & 0 & \nu
 \end{pmatrix},\ \nu\in\mathbb{C} \right\} .
 \]

 Consider now the $({\rm NVE})_1$ and the fundamental matrix solution $\Phi(t)$ at $t=\infty$. In a neighborhood of the irregular singularity $t=\infty$, the matrix $\Phi(t)$ is written
 as
 \[
 \Phi(t)=H_1(t) t^{\Lambda_1} \exp (Q_1 t) ,
 \]
 where $H_1(t)$ is the holomorphic matrix
 \[
 H_1(t)=\begin{pmatrix}
 1 & 0 & -\alpha t^{-2} & -\alpha t^{-2}\\
 0 & 1 & -\alpha t^{-2} & -\alpha t^{-2}\\
 0 & 0 & 1 & 0\\
 0 & 0 & 0 & 1
 \end{pmatrix} .
 \]
 The matrices $\Lambda_1$ and $Q_1$ are given by
 \[
 \Lambda_1={\rm diag} (\alpha-1, \alpha-1, -\alpha, -\alpha),\qquad
 Q_1={\rm diag}(-1, -1, 1, 1) .
 \]
 Since we do not observe non-trivial Stokes phenomena, the local Galois group $G_{\infty}$ is generated topologically by the formal monodromy $\hat{M}_{\infty}$ and the
 exponential torus $\mathcal{T}_{\infty}$. The formal monodromy corresponding to a loop around $t=\infty$ is nothing but $(M_0)^{-1}$, that is, $\hat{M}_{\infty}=(M_0)^{-1}$. Therefore, we can consider
 the local differential Galois group $G_0$ at the origin as a subgroup of the local differential Galois group $G_{\infty}$ at $t=\infty$. Thus the Galois group $G_1$ of the $({\rm NVE})_1$
 coincides with the local differential Galois group $G_{\infty}$ at $t=\infty$. For the exponential torus, we have
 \[
 \mathcal{T}_{\infty}=\begin{pmatrix}
 \lambda^{-1} & 0 & 0 & 0\\
 0 & \lambda^{-1} & 0 & 0\\
 0 & 0 & \lambda & 0\\
 0 & 0 & 0 & \lambda
 \end{pmatrix} ,
 \]
 where $\lambda\in\mathbb{C}^*$.

 As a result, we find that when $\alpha\in\mathbb{Q}$ the connected component $(G_1)^0$ of the differential Galois group $G_1$ coincides with $\mathcal{T}_{\infty}$. When $\alpha\notin\mathbb{Q}$ the group
 $(G_1)^0$ is generated by $\hat{M}_{\infty}$ and $\mathcal{T}_{\infty}$. Summarily, the group $(G_1)^0$ is defined as
 \[
 (G_1)^0=\left\{\begin{pmatrix}
 \mu^{-1} & 0 & 0 & 0\\
 0 & \mu^{-1} & 0 & 0\\
 0 & 0 & \mu & 0\\
 0 & 0 & 0 & \mu
 \end{pmatrix}, \ \mu\in\mathbb{C}^*\right\} ,
 \]
 which is an Abelian group.
 %This ends the proof.
 \end{proof}

\end{document}
